\section{Implementation}
\label{sec:implementation}

To make the design-to-system path explicit, Table~\ref{tab:design-impl-map} summarizes how the two design aspects in Chapter~3 map to implementation focus in this chapter.

\begin{table}[t]
\centering
\caption{Mapping from design aspects to implementation mechanisms.}
\label{tab:design-impl-map}
\begin{tabular}{p{0.38\linewidth} p{0.56\linewidth}}
\hline
\textbf{Design Aspect} & \textbf{Implementation Focus} \\
\hline
Workload-aligned reuse under drift & Bounded hint state, online utility-aware gating, warm-start seed adaptation, and fallback-on-drift safety. \\
Backend-decoupled control interface & Policy/Adapter/Runtime layering, backend-agnostic control signals, and runbook-configurable reproducible execution. \\
\hline
\end{tabular}
\end{table}

\subsection{Implementation for Design Aspect I: Workload-Aligned Reuse under Drift}
To operationalize the first design aspect, reuse is implemented as a guarded online decision process rather than a fixed acceleration path.
The core runtime state is a bounded hint table, where each entry maintains confidence, freshness, and recent utility signals.
All updates are incremental and localized, ensuring that policy maintenance remains lightweight under high-rate streams.

At query time, the policy layer applies utility-aware gating: hints are activated only when recent evidence indicates positive gain; otherwise, the system reverts to native traversal.
For dynamic HNSW backends, accepted hints are translated into warm-start seeds and controlled expansion signals.
When drift weakens confidence, fallback is triggered immediately to prevent stale reuse from degrading recall.
This mechanism instantiates the ``selective and self-correcting'' behavior introduced in Section~3 while preserving throughput in stable phases.

\subsection{Implementation for Design Aspect II: Backend-Decoupled Control Interface}
To operationalize the second design aspect, StreamSeed is integrated into CANDOR-Bench as a pluginized module with explicit separation between policy and backend mechanics.
The implementation is organized into three cooperating layers:
\begin{itemize}
	\item \textbf{Policy layer}: AQHC state machine, hint scoring, gating, and eviction.
	\item \textbf{Adapter layer}: translation between backend-agnostic hints and backend-native index operations.
	\item \textbf{Runtime layer}: runbook-driven scheduling, maintenance orchestration, and statistics logging.
\end{itemize}

This interface-first structure preserves existing experiment entry points while allowing backend-specific optimizations to evolve independently.
Major controls (hint budget, gating thresholds, refresh aggressiveness, and maintenance cadence) are exposed through runbook and algorithm configuration files to ensure reproducibility.
Consequently, the same policy implementation can be compared across event rates, datasets, and dynamic backends under consistent instrumentation.
